% Guia do professor — Capítulo 5: Estabilidade Atmosférica
\section{Capítulo 5 --- Estabilidade Atmosférica}

\subsection*{Visão geral e ritmo}

O capítulo-síntese da primeira metade do curso: Caps.~1--4 se encontram nos
diagramas termodinâmicos e no método da parcela. Empuxo, Brunt--Väisälä,
estabilidade estática (com o método \emph{não local} que é marca registrada
do Stull) e número de Richardson. \textbf{Ritmo sugerido: 2 aulas de 2\,h}
(o capítulo é longo no livro; estas notas destilam o essencial).

\begin{itemize}
  \item \textbf{Aula 1} — diagramas termodinâmicos (construção passo a
        passo no quadro), método da parcela, empuxo, dedução do oscilador
        de Brunt--Väisälä. Fechar com Casos~1--2 (15\,min).
  \item \textbf{Aula 2} — estabilidade local × não local (o coração
        pedagógico da aula), Richardson e turbulência, leitura de perfis
        (tropopausa, camada de mistura, inversões). Casos~3--4 (15\,min).
\end{itemize}

\subsection*{Objetivos de aprendizagem}

\begin{enumerate}
  \item ler e usar um Skew-$T$: adiabáticas secas/úmidas, isoumas, LCL,
        trajetória de parcela;
  \item calcular o empuxo via temperatura virtual e estimar velocidades de
        termais;
  \item deduzir $N^2$ e interpretar o período de Brunt--Väisälä em
        fenômenos visíveis;
  \item aplicar o método não local e explicar por que o critério local
        falha;
  \item calcular $Ri$ e prever turbulência de céu claro;
  \item identificar tropopausa, camada de mistura e inversões num perfil.
\end{enumerate}

\subsection*{Pré-requisitos a reativar}

$T_v$ (Cap.~1), $\theta$ e $\Gamma_d$ (Cap.~3), LCL e $r_s$ (Cap.~4),
oscilador harmônico (física básica — os alunos vão sorrir ao reencontrá-lo
na atmosfera).

\subsection*{Armadilhas conceituais frequentes}

\begin{itemize}
  \item \textbf{``Gradiente ambiental $6{,}5 <$ adiabático $9{,}8$ $\Rightarrow$
        sempre estável''}: verdade para o ar \emph{não saturado} em média —
        mas o método da parcela compara com a \emph{adiabática úmida} acima
        do LCL, e a instabilidade condicional aparece. Preparar bem: é a
        ponte para o Cap.~14.
  \item \textbf{Estabilidade só pelo gradiente local}: a armadilha central;
        o exemplo da tarde convectiva desfaz.
  \item \textbf{Confundir $N$ com frequência de oscilação visível}: nuvens
        em faixa têm espaçamento espacial $U\cdot P_{BV}$; a conversão
        tempo$\to$espaço confunde.
  \item \textbf{$Ri$ como critério simétrico}: a transição
        laminar$\to$turbulento ($Ri<0{,}25$) e a volta ($Ri>1$) não
        coincidem — histerese.
  \item \textbf{Esquecer $T_v$ no empuxo}: em ar úmido tropical o erro
        chega a alguns décimos de K — exatamente a ordem do empuxo de um
        termal.
\end{itemize}

\subsection*{Demonstração de baixo custo}

Garrafa PET com água + azeite (ou água quente colorida embaixo de fria):
oscilações de interface ilustram Brunt--Väisälä; agitar mostra a transição
para turbulência (Richardson qualitativo). Vídeo de ondas de
Kelvin--Helmholtz em nuvens fecha a Aula 2 com aplauso.
\textbf{Material de apoio}: o \S5.11 do livro traz diagramas
termodinâmicos em página inteira (emagrama, Skew-$T$, tefigrama etc.)
prontos para imprimir — distribuir cópias para o trabalho manual em
sala e nas provas com consulta.

\subsection*{Problemas sugeridos do livro (exercícios do Cap.~5)}

Aquecimento: uso do diagrama termodinâmico (plotar sondagem, achar LCL);
cálculo de $N$ e período; classificação de estabilidade de perfis dados.
Aprofundamento: método não local em perfis com inversão; $Ri$ com
cisalhamento realista. Síntese: ``por que o topo da camada de mistura
coincide com a base das nuvens em dias de bom tempo?''.
\emph{Confira a numeração na sua cópia (v1.02b).}

\subsection*{Uso do notebook \texttt{cap05\_estabilidade.ipynb}}

\begin{center}
\begin{tabular}{lll}
\hline
Caso & Conteúdo & Momento sugerido \\
\hline
1 & Diagrama do zero + Skew-$T$ com parcela e área de empuxo & Aula 1, após diagramas \\
2 & Oscilador de Brunt--Väisälä (EDO + amortecimento) & Aula 1, após $N^2$ \\
3 & Estabilidade não local implementada & Aula 2, após o método \\
4 & Perfil de $Ri$ e camadas de CAT & Aula 2, após Richardson \\
\hline
\end{tabular}
\end{center}

O Caso~1 tem intenção dupla: desmistificar o diagrama (construindo-o) e
plantar o CAPE (área sombreada) para o Cap.~14. O Caso~3 é a implementação
fiel do algoritmo do livro — raro em qualquer curso.

\subsection*{Exercícios complementares e soluções}

Nas notas: T1--T5 (empuxo com \emph{liquid loading}, dedução de $N^2$,
períodos por camada, CAT no jato, contraexemplo local × não local) e
N1--N4. Soluções em \texttt{cap05\_solucoes.ipynb} (\emph{não distribuir}):
T2 com \texttt{sympy} (solução do oscilador), N1 com amortecimento, N3 com
o algoritmo comentado linha a linha. Pares: T2$\leftrightarrow$N1,
T5$\leftrightarrow$N3.

\subsection*{Conexões com o resto do curso}

Diagramas + parcela $\to$ CAPE e tempestades (Cap.~14); não local $\to$
camada limite (Cap.~18); inversões $\to$ dispersão de poluentes (Cap.~19);
$N$ $\to$ ondas de montanha (Cap.~17); $Ri$ $\to$ turbulência e aviação.
